\subsection{Marco Teórico}

El presente marco teórico no es un simple listado de definiciones, sino un análisis enfocado en el problema principal de esta investigación: el desorden y la separación de los datos clínicos y administrativos en un consultorio dental pequeño, y cómo un sistema integrado puede ayudar a reducir el tiempo que se pierde en estas tareas. Para organizar las ideas, el texto se divide en tres ejes analíticos que responden directamente a los objetivos específicos del estudio: el diagnóstico de la fragmentación de datos, el uso del modelado de procesos para diseñar una solución y la medición del tiempo para comprobar el impacto real del sistema.
\subsubsection{La fragmentación de la información como un problema de gestión}

En el área de la informática médica, se sabe que los problemas para manejar la información en los consultorios pequeños no son fallas del momento, sino un problema de fondo donde los datos están totalmente divididos. Los estudios especializados en odontología confirman que tener archivos sueltos y la falta de reglas comunes son las barreras más repetidas en el sector \cite{ballester2022current}. Esta falta de orden no solo vuelve lento el trabajo diario, sino que daña la continuidad de la atención, porque trabajar con datos incompletos hace que sea más difícil tomar buenas decisiones sobre los tratamientos de los pacientes \cite{gurupur2024need}.

La consecuencia más clara de este desorden es el aumento del papeleo para el personal. Se ha comprobado que las tareas de oficina exageradas son la fuente principal de cansancio y estrés en los empleados de salud \cite{murad2024measuring}. Pasar datos a mano de un registro a otro quita tiempo que debería usarse para atender a las personas. Además, trabajar con carpetas físicas y hojas sueltas vuelve todo más propenso a equivocaciones, ya que los experimentos demuestran que las personas que usan papel cometen un promedio de 9.3 omisiones de datos importantes por cada paciente en comparación con quienes usan sistemas digitales \cite{pmc2024digital}.

En Ecuador, esta situación se vuelve más delicada por las exigencias de las leyes de salud. El Acuerdo Ministerial 5216-A del Ministerio de Salud Pública exige que los consultorios garanticen el archivo seguro, la confidencialidad y el registro de quién revisa las historias clínicas \cite{msp2016reglamento}. Sin embargo, los análisis técnicos demuestran que la mayoría de los consultorios pequeños tienen fallas graves para cumplir con estas reglas de protección de datos personales \cite{martinez2022proteccion}. A esto se suma que los formularios físicos obligatorios, como el Formulario 033, suelen llenarse mal, y menos del 50\% de los expedientes en papel cuentan con los diagnósticos o planes médicos completos \cite{uce2016caracteristicas}.

Para solucionar esto, la ciencia propone unir la información en un solo sistema. Centralizar los datos ayuda a que la información generada en la consulta pase de forma automática a los cobros y a las recetas, mejorando el servicio y eliminando las tareas repetitivas \cite{ovreiu2024empowering}. En el caso de una clínica dental, ordenar el flujo de trabajo y conectar las herramientas digitales reduce los tiempos de espera y aumenta la productividad del consultorio, evitando que el personal tenga que copiar los mismos datos en diferentes lugares \cite{batheja2025optimizing}.

\subsubsection{El modelado de procesos para el diagnóstico y el diseño del software}

Para construir un software útil, primero se debe entender a fondo cómo trabaja el consultorio hoy en día. La ingeniería de software demuestra que los programas diseñados sin revisar el flujo de trabajo real terminan empeorando los problemas que querían arreglar. Analizar paso a paso las tareas junto con el personal, antes de instalar cualquier tecnología, ayuda a que el sistema se adapte a la realidad del negocio y a que los empleados lo acepten con mayor facilidad \cite{staras2021using}. Cuando los sistemas no se alinean con el camino real que sigue el paciente, aparecen fallas graves en la entrega de recetas y en los controles posteriores a la consulta \cite{makerere2024business}.

Para representar estos flujos con claridad, no bastan las descripciones habladas, sino que se necesitan diagramas estructurados que muestren quién hace cada tarea y dónde se rompe la comunicación \cite{antonacci2021process}. La notación gráfica BPMN es la herramienta más recomendada para esto, aunque en el sector salud plantea retos especiales debido a las variaciones en la atención de cada paciente y a que el doctor y la secretaria trabajan al mismo tiempo con datos compartidos \cite{pufahl2022bpmn}.

Dibujar los procesos actuales sirve como base para identificar qué funciones debe tener el nuevo programa. Estos requisitos de software se pueden organizar según el ciclo de vida de la información, abarcando desde cómo se guardan los datos del paciente hasta las seguridades del sistema \cite{kinast2023functional}. Además, para que el programa maneje las reglas del Formulario 033 de forma automática, la teoría aconseja usar herramientas de gestión de reglas de negocio, logrando que el software valide los códigos de salud sin que el usuario tenga que memorizarlos \cite{nelson2014business}. 

El mapa de procesos permite ver la diferencia entre cómo funciona la clínica hoy y cómo debería funcionar en el futuro con la ayuda del software \cite{hoang2026data}. Este enfoque ayuda a eliminar pasos innecesarios que solo quitan tiempo, asegurando una atención más ágil \cite{alogla2025revisiting}. Finalmente, esto permite diseñar una base de datos limpia y ordenada desde el inicio, superando las típicas fallas de los sistemas dentales viejos que suelen tener archivos mal documentados y desconectados \cite{duncan2020structuring}.

\subsubsection{La medición del tiempo como prueba del éxito del proyecto}

La meta final de ordenar y conectar la información es lograr un impacto positivo en el tiempo que se pierde en los trámites del consultorio. Elegir el tiempo promedio por paciente como la variable para medir el éxito del proyecto es una decisión respaldada por la ciencia, ya que es el indicador más claro para evaluar la carga de trabajo de la oficina \cite{murad2024measuring}. Las pruebas en centros médicos reales confirman que el uso de registros electrónicos reduce de forma directa los minutos dedicados al papeleo por cada paciente atendido \cite{pmc2021electronic}.

Para medir estos minutos de forma exacta en el consultorio, la metodología estándar es el estudio de tiempos y movimientos \cite{cureus2022purpose}. Esta técnica consiste en observar directamente las tareas de la clínica y cronometrar cuánto dura cada evento del ciclo administrativo, como abrir la ficha del paciente, escribir la receta o registrar el pago \cite{lopetegui2014time}. Al usar este método en el campo dental, los investigadores han descubierto que los tiempos varían mucho según el rol de la persona, por lo que es vital medir por separado el trabajo del odontólogo y el de la secretaria \cite{schwei2016exploring}.

La mejor forma de comprobar si el sistema funciona es mediante un diseño de investigación de antes y después \cite{pmc2021electronic}. Al medir los tiempos bajo las mismas condiciones de la clínica antes de poner el programa y volver a medirlos meses después de su instalación, se puede demostrar matemáticamente si la herramienta tecnológica cumplió con su objetivo.

Por último, los autores recuerdan que el éxito de un software no solo depende de que esté bien programado, sino de factores humanos como la usabilidad, la capacitación y la buena disposición del personal para cambiar su forma de trabajar \cite{haddad2017perceived}. En los países en desarrollo, la falta de costumbre tecnológica y el miedo al cambio son los mayores obstáculos para que un sistema funcione \cite{tun2023clinical}. Por esta razón, integrar a los usuarios en las etapas de diseño ayuda a que sientan la herramienta como suya y se alcancen los resultados de tiempo esperados \cite{saladdin2025information}.